\section{Theoretical Background and Neural Network Formulations}

\subsection{Mathematical Principles and Optimization Dynamics}
The core mathematical formulation governing the forward and backward passes of the architecture is defined by:
\begin{equation}
  \mathcal{L}_{\text{CIoU}} = 1 - \text{IoU} + \frac{\rho^2(b, b^{gt})}{c^2} + \alpha v
\label{eq:core_math}
\end{equation}
with multi-scale Feature Pyramid Network lateral fusion defined by $P_k = \text{Conv}_{3\times3}(\text{Conv}_{1\times1}(C_k) + \text{Upsample}(P_{k+1}))$.
In deep architectures with $L$ sequential parametric layers, the composite mapping $f(x) = f_L(f_{L-1}(\dots f_1(x)))$ is optimized via stochastic gradient descent (SGD) or adaptive moment estimation (Adam). The parameter update for weight tensor $W^{(l)}$ at time step $t$ follows:
\begin{equation}
  m_t = \beta_1 m_{t-1} + (1 - \beta_1) g_t, \quad v_t = \beta_2 v_{t-1} + (1 - \beta_2) g_t^2, \quad W_t = W_{t-1} - \frac{\eta}{\sqrt{\hat{v}_t} + \epsilon} \hat{m}_t
\end{equation}
where $g_t = \nabla_W \mathcal{L}$ denotes the empirical gradient, and $\hat{m}_t, \hat{v}_t$ represent bias-corrected first and second moment vectors.

\subsection{Vanishing/Exploding Gradients and Structural Solutions}
In deep convolutional and recurrent architectures, backpropagating gradients through long computational chains risks exponential vanishing or explosion:
\begin{equation}
  \frac{\partial \mathcal{L}}{\partial h_1} = \frac{\partial \mathcal{L}}{\partial h_T} \prod_{k=1}^{T-1} \frac{\partial h_{k+1}}{\partial h_k} = \frac{\partial \mathcal{L}}{\partial h_T} \prod_{k=1}^{T-1} \text{diag}(1 - \tanh^2(a_{k+1})) W^T
\end{equation}
If the spectral radius $\rho(W) < 1$, gradients decay toward zero exponentially as $T$ increases, blinding the network to long-range dependencies. Conversely, if $\rho(W) > 1$, gradients explode. Residual skip connections ($x + \mathcal{F}(x)$) and gated recurrent cells (LSTM constant error carousels) resolve this by providing an additive gradient identity highway $\frac{\partial (x + \mathcal{F})}{\partial x} = I + \frac{\partial \mathcal{F}}{\partial x}$, ensuring sustained gradient backpropagation across hundreds of layers \cite{he2016deep, hochreiter1997long}.

\subsection{Regularization and Generalization Bounds}
To prevent over-parameterized neural networks from memorizing training noise, modern architectural regularization techniques are integrated:
\begin{itemize}
  \item \textbf{Batch Normalization:} Normalizes layer inputs across mini-batches $\hat{x} = \frac{x - \mu_B}{\sqrt{\sigma_B^2 + \epsilon}}$, smoothing the loss landscape and accelerating convergence.
  \item \textbf{Spatial Dropout:} Randomly zeroes entire feature channels during training with probability $p$, breaking co-adaptation between adjacent convolutional kernels.
  \item \textbf{Weight Decay ($L_2$):} Adds a quadratic penalty $\frac{\lambda}{2} \|W\|_2^2$ directly to the loss, penalizing large weights and constraining effective model capacity.
\end{itemize}
